\subsection{Auxiliary Study: Intermediate Feedback on LIBERO-Spatial}\label{sec:libero_results}
This study examines one question in isolation: whether intermediate feedback helps full-task RL. It uses a different backbone and does not run the adaptation schedule. It also does not test the milestone credit rule of Section~\ref{sec:composition}.

We fine-tune an OpenVLA-OFT policy with LoRA on complete instructions for the ten LIBERO-Spatial tasks~\cite{openvlaoft,libero}, using our own implementation of SimpleVLA-RL~\cite{simplevlarl}. The terminal-only row is our run under this study's budget, not a result from the SimpleVLA-RL publication. All runs start from one shared checkpoint with 61.0\% success; their final rates are measured after training. Each endpoint uses 300 validation rollouts from one training seed. The dense-plus-phase and terminal-only runs share terminal verifier coefficient 5 and 160 updates. The phase-only run differs in terminal coefficient (1), phase weight, and budget (110 updates), so we report it as a separate configuration rather than as a matched ablation. Dynamic sampling can make equal update counts consume different numbers of rollouts. Environment interactions and GPU-hours per run were not retained, so we compare runs only by updates.
